\section{Mechanistic Interpretability：内部证据的希望与限制}

Behavioral evaluation 只能观察输入输出，\term{mechanistic interpretability}（机制可解释性）尝试从 activation、feature 和 circuit 理解模型内部计算。Anthropic 的 dictionary-learning 工作把单个 neuron 的混合表示分解为更稀疏的 feature，为“某种概念何时激活、如何影响下游”提供研究工具。

\subsection{Feature、circuit hypothesis 与 intervention}

\term{feature} 是从高维 activation 中提取的可解释模式，不一定对应单个 neuron；\term{circuit} 是多个 feature/attention/MLP interaction 形成的计算路径。研究流程从 activation collection、feature discovery 到 behavior correlation，再用 ablation、steering 或 counterfactual intervention 检验因果。

\lecturefigure{14-interpretability-loop.png}{机制可解释性从 activation 提取 feature，形成 circuit hypothesis，再用 intervention 验证。}{Anthropic monosemanticity research；本地字幕 38:17--39:03；概念重绘。}

\begin{knowledgebox}{读图：可读 feature 不等于完整“读心”}
Feature label 来自激活样本和人类解释，可能遗漏 polysemy、distributed representation 或 context dependency；circuit hypothesis 需要 causal intervention；局部解释也不能证明全局安全。Interpretability 更适合作为 evaluation、red team 和 anomaly investigation 的补充证据，而不是单独 release gate。
\end{knowledgebox}

\begin{warningbox}{不要把 anthropomorphic label 当模型意图}
“欺骗”“权力”“自我保护”等标签可以帮助检索行为模式，但模型内部 feature activation 不自动意味着持久目标或主观意图。报告应说明 operational definition、dataset、false positive、intervention result 和 alternative explanation。
\end{warningbox}

\subsection{本章小结}

Mechanistic interpretability 提供内部 evidence，但仍处于研究早期。Feature discovery、causal validation 与 behavioral evaluation 必须相互校验，不能把一张 activation 图当成安全证明。

